\section{Verification Strategy and Threat Model}
\label{sec:verification}

We follow the libcrux\cite{libcrux} pattern: a portable Rust path
serves as the source of truth for formal verification (via
hax\cite{hax-toolchain} extracting to F$^*$), while a parallel
\texttt{asm}-feature path provides performance.  The two are kept
byte-equivalent through differential testing.  This section
enumerates the verification scope, what each layer catches, and
which residual trust assumptions remain.

\subsection{Layered verification}

\begin{table}[h]
\centering
\small
\begin{tabular}{l l l}
\toprule
\textbf{Layer} & \textbf{Tool} & \textbf{What it catches} \\
\midrule
Algorithmic CT       & hax + F$^*$ + \texttt{secret\_integers}
                     & Branches, indexing, division on secrets in Rust \\
Memory safety        & F$^*$ type-check                 & Panics, integer overflow, OOB access \\
Spec equivalence     & F$^*$ refinement proofs         & Functional correctness vs.\ Algorithm~4.x \\
Compiled-code CT     & LLVM CT-coverage runs           & Compiler not reintroducing branches \\
Empirical CT         & \texttt{dudect}, \texttt{tacet} & Wall-clock timing variance under load \\
Asm-vs-Rust          & Differential property tests     & Output equivalence on N random inputs \\
KAT correctness      & 100 NIST PQC vectors            & End-to-end interop with C reference \\
\bottomrule
\end{tabular}
\caption{Verification layers and the property each provides.}
\label{tab:verification-layers}
\end{table}

\subsection{What hax verification provides}

A successful hax extraction plus F$^*$ type-check on the portable
Rust path proves:
\begin{itemize}[nosep]
  \item \textbf{Panic freedom.}  No path through the verified
        modules can \texttt{unwrap}, panic on integer overflow,
        index out of bounds, or trigger a \texttt{debug\_assert!}.
  \item \textbf{Memory safety by construction.}  Hax refuses to
        extract \texttt{unsafe}, raw pointers, FFI, or inline asm.
        The verified path is structurally safe.
  \item \textbf{No computation over secret data.}  With one
        annotation per secret-bearing field
        (\texttt{Coordinate}, \texttt{BigInt} representations of
        secret keys, etc.), F$^*$ rejects any branch on, indexing
        by, or division/modulo by a secret value.  Existing
        \texttt{subtle::Choice} / \texttt{ConditionallySelectable}
        usage flips from convention to proof.
\end{itemize}

The codebase's \texttt{TODO(ct)} markers (one per known
variable-time path on secret data, see \S\ref{sec:ct}) become the
exact set of locations where the \texttt{secret\_integers}
annotation must succeed; if any remain after CT hardening, F$^*$
extraction fails and surfaces the gap.

\subsection{Limitations of hax-based verification}

Hax verifies the Rust source as written; the property landscape
\emph{below} that level is not in scope.  We list the limitations
explicitly so users understand the residual trust:

\begin{itemize}[nosep]
  \item \textbf{Cache-timing channels} on data-cache or
        instruction-cache patterns that depend on secret-derived
        prior state.  Hax models direct indexing
        ($\texttt{arr}[\texttt{secret\_idx}]$) but not prefetcher
        behavior, code-cache occupancy of differently-sized
        callees, or set-associative cache contention with
        co-tenants.
  \item \textbf{Variable-time hardware instructions} ($\texttt{DIV}$
        on \texttt{x86\_64}, $\texttt{UDIV}$ on \texttt{aarch64},
        \texttt{LZCNT} / \texttt{POPCNT} on some older parts).
        F$^*$ rejects \emph{secret-dependent} division explicitly,
        but a non-rejected ordinary integer divide still has data-dependent
        latency on the underlying hardware.
  \item \textbf{Speculative execution} (Spectre v1/v2/v4, RIDL,
        ZenBleed, \dots).  Hax models architectural semantics, not
        the speculative paths the CPU explores.  Mitigation is
        out-of-scope for the verified path; users requiring
        Spectre resistance must add platform-specific barriers.
  \item \textbf{Power and EM side channels.}  Outside any
        compiler-level analysis.
  \item \textbf{Allocator timing.}  A first-time
        \texttt{Vec::with\_capacity} that triggers a syscall is a
        side channel hax does not see.  Hot-path allocations on
        secret-derived sizes are addressed by code review and
        by structural use of fixed-width \texttt{BigInt<N>}.
  \item \textbf{Cross-process timing attacks.}  Hyperthreading,
        SMT contention, scheduler-induced bandwidth.  Mitigations
        are at the OS/sysadmin layer.
  \item \textbf{Asm path.}  Anything inside an \texttt{asm!} block
        or behind \texttt{cfg(feature = "asm")} is unverified by
        hax.  The asm path is treated as
        \emph{trusted-but-tested}: differential property tests
        prove byte-equivalence against the (verified) portable
        path, but not independent CT on the asm itself.
\end{itemize}

\subsection{Residual trust assumptions}

A user of \texttt{sqisign-selkie} with both hax verification and
the asm path enabled is trusting:
\begin{enumerate}[nosep]
  \item Correctness of \texttt{rustc}, LLVM, and the F$^*$ type
        checker.
  \item The hax extraction toolchain as a translation between
        Rust and F$^*$.
  \item The C reference implementation's algorithm choices, where
        we follow them (e.g., dual-LLL ball sampling, the C ref's
        kernel construction in
        \texttt{compute\_dim2\_isogeny\_challenge}).
  \item Hardware implementation of the chosen targets meeting the
        side-channel claims (variable-time instructions absent on
        the asm hot path; no Spectre exploitation in the
        deployment context).
  \item The asm path's correctness (byte-equivalent to the
        verified Rust path for all encountered inputs, beyond the
        $N$ random inputs covered by differential tests).
\end{enumerate}

\subsection{Roadmap and current status}

The current state of each layer.  All items below are explicitly
\textbf{not yet shipped} unless marked \textbf{done}; the roadmap is
intentionally honest about what verification claims this codebase
\emph{does} and \emph{does not} support today.

\begin{itemize}[nosep]
  \item KAT correctness (Layer~7): \textbf{done}.
        \emph{Verify} accepts the C reference's 100 NIST-I KAT
        signatures byte-for-byte; \emph{sign} produces
        byte-identical signatures on all 100 KAT seeds and on the
        full 103-vector keygen suite.  All signatures cross-verify
        against the C reference harness.
  \item Empirical CT (Layer~5): \textbf{infrastructure present;
        not a verification claim}.  \texttt{dudect} and
        \texttt{tacet} are wired into CI with baseline measurements,
        but the codebase is not yet constant-time on secret-derived
        paths (see~\S\ref{sec:ct}).
  \item Asm-vs-Rust differential testing (Layer~6):
        \textbf{NOT STARTED}.  Activates when the asm path lands.
  \item Algorithmic CT via hax + \texttt{secret\_integers}
        (Layer~1): \textbf{NOT STARTED}.  Hax extraction is queued
        behind asm-Fp to avoid re-verifying an evolving impl.
  \item F$^*$ memory safety / panic freedom (Layer~2):
        \textbf{NOT STARTED}.  Lands together with hax extraction.
  \item Spec equivalence proofs (Layer~3): \textbf{NOT STARTED}.
        Out of scope for the v0.1 tag; deferred indefinitely.
  \item Asm formal verification: \textbf{NOT STARTED}.  See the
        open work item below.
\end{itemize}

\paragraph{Open work item: formal verification of the asm path
(\emph{TODO, not yet started}).}
Hax does not see asm.  The trust gap above (item~5 in ``residual
trust'') is the standard libcrux model: byte-equivalent diff tests
against a verified Rust source.  For higher-assurance deployments,
a separate effort using EverCrypt/Vale-style asm
verification~\cite{vale-asm} would close that gap with a formal
proof that the asm respects the same functional + CT specification
F$^*$ verifies on the Rust path.  This is tracked here as a
TODO for a post-tag milestone; shipping it requires writing Vale
specifications for the Fp/Fp$^2$ hot paths and integrating an
additional verification pipeline.  The plan, in order, is:

\begin{enumerate}[nosep]
  \item \textbf{(in progress)} Tag spec/C-ref correctness (current
        milestone target).
  \item \textbf{(not started, queued next)} Implement asm-Fp on
        \texttt{aarch64} and \texttt{x86\_64}; diff-test against
        portable Rust.
  \item \textbf{(not started)} Hax-extract the portable Rust path
        to F$^*$; verify memory safety, panic freedom, and
        \texttt{secret\_integers} CT.
  \item \textbf{(not started, future)} Vale verification of the
        asm path against the same F$^*$ specification.
\end{enumerate}

Until each step lands, the corresponding verification claim
\emph{does not apply} and should not be relied on.  Users should
treat this section as a forward-looking roadmap, not a statement
of current properties.

\subsection{What hax cannot replace}

Independent of the side-channel layers above, hax is not a
substitute for:

\begin{itemize}[nosep]
  \item Algorithmic correctness review.  F$^*$ proves the code
        matches a \emph{spec}; the spec itself can be wrong, and
        SQIsign's specification has historically been ambiguous in
        ways our bug catalog (\S\ref{sec:bugs}) documents.
  \item KAT testing.  F$^*$ proves correctness up to the
        specification's expressed properties; KAT vectors prove
        interoperability with deployed implementations.  Both are
        needed.
  \item Negative testing.  Wycheproof-style adversarial vectors
        catch failure modes (malleable signatures, parser
        confusion) that functional verification often does not.
\end{itemize}
